\begin{trapbox}
\begin{enumerate}
  \item \textbf{混淆任意点公式与短轴端点公式。}
  \[
  \underbrace{\tan\frac\theta2=\frac{c|y_0|}{b^2}}_{\text{椭圆上任意点}}
  \qquad\ne\qquad
  \underbrace{\frac c{|y_0|}}_{\text{一般不成立}}.
  \]
  只有 $|y_0|=b$ 时，二者同为 $c/b$。

  \item \textbf{误把 $y_0$ 当作有符号纵坐标。}
  $\theta$ 是非负角，椭圆上下两半关于 $x$ 轴对称；
  必须使用 $|y_0|$。例如 $P=(0,-b)$ 也取得最大角。

  \item \textbf{误以为焦角确定就唯一确定 $P$。}
  同一椭圆上确实有
  \[
    \theta\text{ 相同}\iff
    S_{\triangle PF_1F_2}\text{ 相同}\iff
    |y_0|\text{ 相同},
  \]
  但它只决定两个坐标的\textbf{绝对值}。
  $0<\theta<\theta_{\max}$ 时通常有四个对称点，
  $\theta=0$、$\theta=\theta_{\max}$ 时分别只有两个点。
  更不能把 ``一个角决定面积'' 推广到\textbf{任意三角形}；
  这里还依赖固定底边 $2c$ 与固定两腰和 $2a$。

  \item \textbf{混淆最小角与最大角的取值位置。}
  $|y_0|=0$ 对应\textbf{长轴端点} $P=(\pm a,0)$，
  此时面积为零，\(\theta_{\min}=0^\circ\)；
  $|y_0|=b$ 对应\textbf{短轴端点} $P=(0,\pm b)$，
  此时面积和焦角同时最大。焦角零度是允许的退化情形，
  不能因为三角形退化就把最小值排除掉。

  \item \textbf{滥用正切和角公式。}
  焦点三角形的三个\textbf{内角}满足
  $\theta+\angle PF_1F_2+\angle PF_2F_1=180^\circ$，
  不是 $\theta=\angle PF_1F_2+\angle PF_2F_1$。
  此外，$\tan\theta$ 在 $\theta=90^\circ$ 时无定义，
  且在钝角区间为负；不能靠 $\tan\theta$ 在全区间判断角度的增减。
  应优先使用 $\tan(\theta/2)$。

  \item \textbf{余弦定理漏掉夹角的对边。}
  $\theta$ 位于 $P$，对边是 $F_1F_2=2c$，故
  $\cos\theta=(r_1^2+r_2^2-4c^2)/(2r_1r_2)$。
  不要误把对边写成 $2a$ 或漏掉平方。

  \item \textbf{从椭圆方程消去 $x_0$ 的代数错误。}
  正确的是
  \[
  x_0^2=a^2\left(1-\frac{y_0^2}{b^2}\right),\qquad
  c^2-e^2x_0^2=\frac{c^2y_0^2}{b^2}.
  \]
  原证明的加减号及化简不一致，应以此式核对。

  \item \textbf{坐标轴方向决定使用哪个坐标。}
  若长轴在 $y$ 轴，即
  $\frac{x^2}{b^2}+\frac{y^2}{a^2}=1$，焦点为 $(0,\pm c)$，
  则高是点到 $y$ 轴的距离 $|x_0|$，对应公式为
  \[
  \boxed{\tan\frac\theta2=\frac{c|x_0|}{b^2}}.
  \]
  更一般地，要取点到\textbf{焦点连线所在直线}的距离 $h$：
  $\tan(\theta/2)=ch/b^2$。

  \item \textbf{取值端点与离心率。}
  对本题 $0<e<1$，最大角严格小于 $180^\circ$；
  $e\to1^-$ 时才有 $\theta_{\max}\to180^\circ$，
  不能把极限值当作某个非退化椭圆实际取到的角。
  当 $e\to0^+$ 时，两个焦点趋于重合，最大角趋于 $0$。
\end{enumerate}
\end{trapbox}
